%% kapa: exam=1 round=31 year=2020 period=1 form=A subject=economics subject_ko=경제학원론 seq=27 number=67
%% source: Q-Net 감정평가사 기출문제 (cst00309) artlSeq=5207219 file=extracted/1008/5207219/A형/A형/2020년도 제31회 감정평가사 1차 시험_1교시_A형.pdf
%% conversion: kordoc 4.24.51 -> tools/gen1.py (auto); review=proofread-sonnet; flags=none
\begin{question}{67}
유동성함정(liquidity trap)에 관한 설명으로 옳은 것을 모두 고른 것은?
\begin{framedtext}
ㄱ. IS 곡선이 수직선이다. \par
ㄴ. LM 곡선이 수평선이다. \par
ㄷ. 재정정책이 국민소득에 영향을 주지 않는다. \par
ㄹ. 화폐수요의 이자율 탄력성이 무한대일 때 나타난다. \par
\end{framedtext}
\begin{choices}
  \item ㄱ, ㄷ
  \item ㄴ, ㄹ
  \item ㄷ, ㄹ
  \item ㄱ, ㄴ, ㄷ
  \item ㄴ, ㄷ, ㄹ
\end{choices}
\end{question}
\answer{2}
